\Askhsh{34026_SOLUTION}{1}
\begin{ylh}{1}
    \item [Ύλη από εικόνα]
\end{ylh}
\begin{alist}
    \item []Δίνεται η συνάρτηση $f(x)=(1-2\ln x)e^{-\frac{x}{x}}$ με $x>0$.
    \begin{rlist}
        \item Να αποδείξετε ότι η συνάρτηση $f$ είναι συνεχής ως άθροισμα συνεχών συναρτήσεων και παραγωγίσιμη με $f'(x)=(1-2\ln x)e^{-\frac{2}{x}}$. Άρα η συνάρτηση είναι γνησίως φθίνουσα στο $(0, +\infty)$. \monades{9}
        \item Να είναι $f(x)=0 \Leftrightarrow 1-2\ln x = 0 \Leftrightarrow \ln x = \frac{1}{2} \Leftrightarrow x=e^{\frac{1}{2}}$. \\
            Από α)i. έχουμε $f \searrow$ για $x>0$, οπότε $x<e^{\frac{1}{2}} \Rightarrow f(x)>f\left(e^{\frac{1}{2}}\right) \Rightarrow f(x)>0$ και $x>e^{\frac{1}{2}} \Rightarrow f(x)<f\left(e^{\frac{1}{2}}\right) \Rightarrow f(x)<0$. \monades{8}
    \end{rlist}
    \item []
    \begin{rlist}
        \item i. Στο διάστημα $(0, +\infty)$ η συνάρτηση $g$ είναι συνεχής ως πηλίκο συνεχών συναρτήσεων και παραγωγίσιμη με $g'(x)=\frac{\frac{1}{x}x^2-(\ln x)2x}{(x^2)^2}=\frac{x-2x\ln x}{x^4}=\frac{x(1-2\ln x)}{x^4}=\frac{1-2\ln x}{x^3}$. \monades{9}
        \item ii. Οπότε για την συνεχή συνάρτηση $g$ συμπληρώνουμε τον πίνακα
        \[
        \begin{tabular}{|c|c|c|c|c|}
            \hline
            $x$ & \cellcolor{gray!20}$0$ & \multicolumn{2}{c|}{$e^{1/2}$} & $+\infty$ \\
            \hline
            $g'(x)$ & \cellcolor{gray!20}$\|$ & $+$ & $0$ & $-$ \\
            \hline
            $g(x)$ & \cellcolor{gray!20}$\|$ & $\nearrow$ & $\frac{1}{2e}$ & $\searrow$ \\
            \hline
        \end{tabular}
        \]
        Άρα η τιμή $g\left(e^{\frac{1}{2}}\right)=\frac{\ln e^{\frac{1}{2}}}{\left(e^{\frac{1}{2}}\right)^2}=\frac{\frac{1}{2}}{e}=\frac{1}{2e}$ είναι η μέγιστη της συνάρτησης $g$. \monades{9}
        \item iii. Η συνάρτηση $g$ είναι συνεχής στο πεδίο ορισμού της $(0, +\infty)$.
        Έχουμε $\lim_{x \to 0^+} g(x) = \lim_{x \to 0^+} \frac{\ln x}{x^2} = \lim_{x \to 0^+} \frac{-\infty}{0^+} = -\infty$ και
        \[
        \lim_{x \to +\infty} g(x) = \lim_{x \to +\infty} \frac{\ln x}{x^2} \xrightarrow{DLP} \lim_{x \to +\infty} \frac{\frac{1}{x}}{2x} = \lim_{x \to +\infty} \frac{1}{2x^2} = 0.
        \]
        Επομένως και λόγω του β)i ερωτήματος το σύνολο τιμών της $g$ είναι το διάστημα $\left(-\infty, \frac{1}{2e}\right)$. \monades{10}
    \end{rlist}
\end{alist}